\documentclass[11pt]{article}
\usepackage[margin=1in]{geometry}
\usepackage[T1]{fontenc}
\usepackage{lmodern}
\usepackage{microtype}
\usepackage{amsmath,amssymb,amsfonts}
\usepackage{booktabs}
\usepackage{array}
\usepackage{enumitem}
\usepackage{graphicx}
\usepackage{placeins}
\usepackage{xcolor}
\usepackage{hyperref}
\usepackage[numbers,sort&compress]{natbib}

\hypersetup{
  colorlinks=true,
  linkcolor=blue!55!black,
  citecolor=blue!55!black,
  urlcolor=blue!55!black
}

\setlist[itemize]{leftmargin=1.6em}
\setlist[enumerate]{leftmargin=1.8em}
\newcommand{\code}[1]{\texttt{#1}}
\newcommand{\R}{\mathbb{R}}
\newcommand{\E}{\mathcal{E}}
\newcommand{\tr}{\operatorname{tr}}
\newcommand{\diag}{\operatorname{diag}}
\newcommand{\vol}{\operatorname{vol}}
\newcommand{\argmin}{\operatorname*{arg\,min}}

\newenvironment{reviewbox}[1]
{\par\medskip\noindent\textbf{#1.}\ }
{\par\medskip}

\newcommand{\paperfigure}[4]{%
\begin{figure}[htbp]
\centering
\includegraphics[width=0.82\linewidth]{#1}
\caption{#2 Reproduced from #3, internal review only. #4}
\end{figure}
}


\title{Paper Review: Convergence of Graph Laplacian with kNN Self-tuned Kernels}
\author{Writer-P10 polished review}
\date{Built 2026-05-15}

\begin{document}
\maketitle

\begin{abstract}
\citet{cheng2022knn} gives a theory-backed account of a practical local-scale
kernel construction: choose a bandwidth at each data point from a
\(k\)-nearest-neighbor radius, use the product of endpoint bandwidths in a
radial kernel, and ask which continuum Laplacian the graph approximates. The
paper proves uniform relative consistency of the kNN scale
\(\widehat\rho\) for \(\bar\rho=p^{-1/d}\), then uses that result to prove
Dirichlet-form and pointwise graph-Laplacian convergence for a family of
self-tuned kernels indexed by \(\alpha\). Its main relevance for a
conductance/kernel graph Laplacian review is as a rigorous adaptive
local-scale comparator: it explains how self-tuning changes density bias,
variance in low-density regions, and normalization targets. It should not be
read as a description of the current \code{fit.rdgraph.regression()} smoother,
whose default conductance is overlap-density/Riemannian-complex based rather
than a kNN self-tuned kernel conductance.
\end{abstract}

\section{Why This Paper Matters}

The self-tuned kernel idea is simple enough to look like a heuristic. If a
point lies in a sparse part of the sample, enlarge its bandwidth; if it lies
in a dense part, shrink it. In the Zelnik-Manor and Perona construction
\citep{zelnikmanor2004selftuning}, the bandwidth at \(x_i\) is the distance
\(\widehat R_i\) from \(x_i\) to its \(k\)-th nearest neighbor, and the
affinity is roughly
\[
  W_{ij}=k_0\!\left(\frac{\|x_i-x_j\|^2}{\widehat R_i\widehat R_j}\right).
\]
This is attractive in high-dimensional data analysis because a single global
bandwidth is rarely appropriate across dense and sparse regions. The
mathematical difficulty is that the kNN radius is random, nonsmooth, and not
derivative-consistent. A standard variable-bandwidth kernel proof cannot
simply assume a smooth bandwidth function and proceed.

Cheng and Wu's contribution is to show that this practical construction still
has a clean continuum interpretation. They prove that the normalized kNN
radius is uniformly close in relative error to \(p^{-1/d}\), where \(p\) is
the sampling density on a compact smooth \(d\)-dimensional manifold. They
then condition on this high-probability event and prove convergence of graph
Dirichlet forms, modified random-walk graph Laplacians, and unnormalized
graph Laplacians. The paper therefore occupies a useful middle ground between
early self-tuning practice and smooth variable-bandwidth diffusion theory:
the local scale is actually estimated from samples, not granted as an oracle.

For SIMODS and \code{gflow}, the paper is important because it supplies a
principled local-scale kernel comparator. It supports testing an adaptive
kernel-weighted graph energy against the existing smoother, especially when
sample density varies strongly. It does not say that the existing
\code{fit.rdgraph.regression()} path is already implementing such a method.
In the current precomputed-graph setting, \code{weight.list} supplies
positive edge lengths used for local ordering and truncation, while the
default smoothing conductance is overlap-density/Riemannian-complex based,
approximately \(c_e^\rho=1/\max(\rho_1(e),10^{-10})\) in the relevant regime.
P10 belongs on the comparator side of the design ledger, not as a
reinterpretation of that current overlap-density behavior.

\section{Problem Setting and Reader Orientation}

The paper assumes samples from a smooth compact boundaryless manifold
\(M\subset\R^D\), with smooth density \(p\) bounded above and below. A reader
should keep three graph objects separate.

First, there is the fixed-bandwidth affinity
\[
  W_{ij}=k_0\!\left(\frac{\|x_i-x_j\|^2}{\epsilon}\right),
\]
whose continuum limits were already well studied. Second, there is the
original self-tuned affinity, where the denominator is
\(\widehat R_i\widehat R_j\). Third, Cheng and Wu introduce a family that
separates the distance scale from an amplitude normalization:
\[
  W_{ij}^{(\alpha)}
  =
  k_0\!\left(
    \frac{\|x_i-x_j\|^2}
         {\epsilon\,\widehat\rho(x_i)\widehat\rho(x_j)}
  \right)
  \frac{1}{\widehat\rho(x_i)^\alpha\widehat\rho(x_j)^\alpha}.
\]
The exponent \(\alpha\) controls which density-weighted operator is targeted.
The product
\(\epsilon\,\widehat\rho(x_i)\widehat\rho(x_j)\) sets the local kernel scale,
while the factor
\(\widehat\rho(x_i)^{-\alpha}\widehat\rho(x_j)^{-\alpha}\) changes the
measure in the limiting Dirichlet form.

The paper deliberately distinguishes the theoretical scale
\(\widehat\rho\) from the raw radius used in the algorithm. The normalized
estimator is
\[
  \widehat\rho(x)
  =
  \widehat R(x)
  \left(\frac{1}{m_0[h]}\frac{k}{N_y}\right)^{-1/d},
  \qquad
  \widehat R(x)
  =
  \inf\left\{r>0:
    \sum_{j=1}^{N_y}\mathbf{1}_{\{\|y_j-x\|<r\}}\ge k
  \right\}.
\]
The stand-alone sample \(Y\) is used to estimate \(\widehat\rho\), while the
graph is built on a sample \(X\). This split simplifies the proof by removing
some dependence among graph weights. Algorithm 1 uses the raw distances
\(\widehat R_i\) and a dimensionless scale \(\sigma_0\), so it does not
require the practitioner to know \(d\):
\[
  W_{ij}
  =
  k_0\!\left(\frac{\|x_i-x_j\|^2}{\sigma_0^2\widehat R_i\widehat R_j}\right)
  \frac{1}{\widehat R_i^\alpha\widehat R_j^\alpha}.
\]
This is the cleanest practical analogue for a SIMODS comparator.

One important terminology caution is that P10 uses kNN distances to estimate
bandwidths. That is not the same thing as proving a kNN graph-support rule.
Edges are determined by the kernel support, by dense implementation, or by
later sparsification. The kNN step supplies local scale; it does not by itself
define the support of the graph analyzed in the continuum theorem.

\section{The kNN Scale Estimate}

The core statistical fact is the uniform relative consistency of
\(\widehat\rho\) for
\[
  \bar\rho=p^{-1/d}.
\]
At a fixed point, Proposition 2.2 gives the relative error as a bias term plus
a variance term,
\[
  \frac{|\widehat\rho(x)-\bar\rho(x)|}{\bar\rho(x)}
  =
  O_p\!\left((k/N_y)^{2/d}\right)
  +
  O\!\left(\sqrt{\frac{\log N_y}{k}}\right),
\]
with constants uniform in \(x\). Lemma 2.1 then supplies the geometric
regularity needed to pass from fixed-point control to uniform control:
\(\widehat R\) is globally Lipschitz, and is smooth away from a finite union
of hyperplanes where nearest-neighbor orderings change. Theorem 2.3 combines
these facts to prove
\[
  \sup_{x\in M}
  \frac{|\widehat\rho(x)-\bar\rho(x)|}{\bar\rho(x)}
  =
  O_p\!\left((k/N_y)^{2/d}\right)
  +
  O\!\left(\sqrt{\frac{\log N_y}{k}}\right)
\]
with high probability.

The contrast with fixed-bandwidth kernel density estimation is one of the
paper's most useful messages. A KDE estimate of \(p\) has relative variance
that worsens like \(p(x)^{-1/2}\) in low-density regions. The kNN radius
estimate instead expands the neighborhood until it contains \(k\) points, so
its relative variance bound is uniform in \(x\), with no corresponding
low-density blowup. That is the mathematical reason self-tuning is useful as
more than a numerical convenience.

\paperfigure
  {../../paper_figures/P10_F01_F02_page15.png}
  {The first page of numerical evidence compares kNN estimation of
  \(\bar\rho=p^{-1/d}\) with fixed-bandwidth KDE estimation of \(p\), and
  sets up the one-dimensional synthetic manifold used later.}
  {Cheng and Wu (2022), Figures 1--2}
  {The top-row relative-error surface in Figure 1 is the key visual evidence:
  the kNN scale estimate has much more spatially uniform relative error than
  the KDE estimate, whose error concentrates in the low-density region.}

There is a subtle but crucial warning in Section 2.3. Although
\(\widehat\rho\) is uniformly \(C^0\)-consistent, its derivative is not
consistent. At differentiability points in the ambient space,
\(|\bar\nabla \widehat R|=1\), so
\[
  |\bar\nabla \widehat\rho(x)|
  =
  \left(\frac{1}{m_0[h]}\frac{k}{N_y}\right)^{-1/d},
\]
which diverges as \(k/N_y\to 0\). Consequently, the kNN scale is not
\(C^1\)-consistent on the manifold, and the paper cannot simply import a
smooth variable-bandwidth proof. The proof strategy is instead to use relative
\(C^0\) control to replace \(\widehat\rho\) by \(\bar\rho\) in carefully chosen
integral and graph quantities.

\section{Limiting Operators and Normalization}

For a smooth positive bandwidth \(\rho\), Cheng and Wu define the differential
operator
\[
  L_\rho^{(\alpha)}
  =
  \Delta
  +2\frac{\nabla p}{p}\cdot\nabla
  +(d-2\alpha+2)\frac{\nabla\rho}{\rho}\cdot\nabla .
\]
Substituting the population scale \(\bar\rho=p^{-1/d}\) gives
\[
  L^{(\alpha)}
  =
  \Delta
  +
  \left(1+\frac{2(\alpha-1)}{d}\right)
  \frac{\nabla p}{p}\cdot\nabla .
\]
Thus \(\alpha\) is not a harmless tuning knob. It determines the density drift
in the continuum limit. In the Dirichlet-form language, the target measure has
density
\[
  p_\alpha=p^{1+2(\alpha-1)/d}.
\]
The special cases are especially useful for comparator design. When
\(\alpha=1\), \(p_\alpha=p\), and the limiting operator is the weighted
Laplacian \(\Delta_p\). When \(\alpha=0\), the construction is closest to the
original unnormalized self-tuned affinity and targets a modified density
\(p^{1-2/d}\). When \(\alpha=1-d/2\), \(p_\alpha\) is constant, so the
Dirichlet-form target and the modified random-walk limiting operator recover
the Laplace--Beltrami geometry. The qualifier matters: the unnormalized
pointwise operator in Theorem 3.6 includes an additional density prefactor
\[
  p^{2(\alpha-1)/d}L^{(\alpha)}f.
\]
Therefore it is imprecise to say that every graph operator recovers
\(\Delta_M\) at \(\alpha=1-d/2\) without also naming the normalization.

The paper analyzes two graph-Laplacian actions. The unnormalized operator is
the scaled sample average associated with \(D-W\). The modified random-walk
operator is a row-normalized average with an additional
\(\widehat\rho(x)^{-2}\) factor; in matrix terms it differs from the usual
\(I-D^{-1}W\) random-walk Laplacian by this local-scale diagonal correction
and by constants/sign. This distinction is relevant for \code{gflow}: a
smoothing penalty \(f^T L f\) is closer in spirit to the Dirichlet-form result
than to a pointwise random-walk generator, while a diffusion or embedding
comparator would need to specify exactly which normalization is being used.

\section{Dirichlet Forms Versus Pointwise Operators}

The cleanest convergence theorem for smoothing is the graph Dirichlet-form
result. The graph form is
\[
  E_N(f,f)
  =
  \frac{1}{\epsilon m_2 N^2}
  \sum_{i,j}
  \epsilon^{-d/2}W_{ij}(f_i-f_j)^2,
\]
up to the paper's equivalent \(f^T(D-W)f\) normalization. Theorem 3.3 shows
that, if \(\epsilon\to 0\), \(N\epsilon^{d/2}\) is large enough, and
\(\epsilon_\rho\to 0\), then
\[
  E_N(f,f)
  =
  E_{p_\alpha}(f,f)(1+O(\epsilon_\rho))
  +O(\epsilon)
  +\hbox{sampling error}.
\]
The important point is that the error from estimating \(\rho\) enters as
\(O(\epsilon_\rho)\). The quadratic form absorbs the roughness of the kNN
scale better than the pointwise operator does.

Pointwise operator convergence is stricter. The modified random-walk theorem
and the unnormalized theorem require \(\epsilon_\rho=o(\epsilon)\), and their
errors include an \(\epsilon_\rho/\epsilon\) term. This is not a minor proof
artifact; it reflects the fact that pointwise graph Laplacian estimates are
sensitive to local oscillation. The paper's weak convergence theorem for the
unnormalized operator removes this \(\epsilon_\rho/\epsilon\) penalty by
testing against a smooth \(\phi\), returning to a form-like statement.

\paperfigure
  {../../paper_figures/P10_F03_page16.png}
  {The first convergence experiment compares errors in the Dirichlet form and
  in pointwise estimates \(L_Nf\) as the kernel scale varies.}
  {Cheng and Wu (2022), Figure 3}
  {For internal synthesis, this figure is most useful because it makes the
  theorem hierarchy visible: the Dirichlet form is less volatile at small
  bandwidths, while pointwise estimates show the expected variance-dominated
  regime.}

This distinction should guide how the paper is reused. If SIMODS is comparing
graph smoothness penalties, Theorem 3.3 is the closest theoretical analogue.
If SIMODS is comparing diffusion operators, embeddings, or local residual
operators, Theorems 3.5 and 3.6 are more relevant, but they demand stronger
scale-estimation accuracy and more careful normalization choices.

\section{Density Robustness and Fixed-Bandwidth Comparators}

The fixed-bandwidth comparison is not just a sidebar. Cheng and Wu compare
their self-tuned family with fixed-bandwidth kernels normalized by an estimated
density \(\widehat p^\beta\). The fixed-bandwidth pointwise variance term has
a factor \(p(x)^{-1/2}\). The self-tuned variance term has the favorable
factor \(p(x)^{1/d}\) in the corresponding theorem. In words: low-density
regions are exactly where fixed bandwidths are most fragile, and exactly
where self-tuned bandwidths expand.

This density robustness is the main reason P10 belongs in a SIMODS/gflow
comparator suite. A global RBF kernel comparator can be informative, but it
has a known failure mode when the graph samples sparse and dense regions
simultaneously. A P10-style comparator would use local radii to reduce that
failure mode. The relevant design degrees of freedom are \(k\), \(\sigma_0\),
\(\alpha\), the choice between dense kernel support and sparsification, and
the choice between Dirichlet-form and random-walk normalization. Reporting all
five choices is essential, because they determine both the finite-sample graph
and the continuum operator being approximated.

The paper's experiments support this reading. Figures 4--6 show the practical
tradeoff among \(\epsilon\), \(k_y\), and whether the radius is estimated from
a stand-alone \(Y\) or from \(X\) itself. A large stand-alone \(Y\) can reduce
bandwidth-estimation error, but with a limited total sample budget, splitting
data can worsen performance by reducing the graph sample size. This is useful
for implementation planning: the stand-alone split is a proof device and a
possible large-data strategy, not a universal prescription.

\section{Embedding Evidence and Its Limits}

The paper also demonstrates the self-tuned kernel on eigenfunction recovery
for \(S^1\) and on MNIST digit embeddings. These figures are useful but should
be cited with discipline. They are empirical evidence that self-tuned kernels
can stabilize embeddings across local density variation; they are not a
spectral-convergence theorem.

\paperfigure
  {../../paper_figures/P10_F08_F09_page19.png}
  {The MNIST experiment compares fixed-bandwidth diffusion-map embeddings with
  two self-tuned kernels and identifies fixed-bandwidth outliers as samples
  with large kNN radii.}
  {Cheng and Wu (2022), Figures 8--9}
  {This screenshot is useful as secondary evidence for low-density
  stabilization. The controlled convergence evidence remains the theorem
  comparison and the synthetic experiments.}

Figure 8 shows that fixed-bandwidth diffusion-map embeddings can be unstable
at smaller scales, while self-tuned kernels remain informative over a wider
range of \(\sigma_0\). Figure 9 identifies the fixed-bandwidth outliers as
points with large seventh-nearest-neighbor distances. This aligns with the
theory's variance story: fixed bandwidths struggle with relatively isolated
points, while adaptive local scales compensate. Still, the paper does not
prove spectral convergence of those MNIST embeddings, nor does it prove that
every self-tuned graph embedding is robust.

\section{Implications for SIMODS and \code{gflow}}

P10 should be used as a source for a new local-scale kernel comparator with a
clear continuum target. A faithful comparator would compute a local radius
\(\widehat R_i\) from graph or ambient distances, choose a dimensionless
\(\sigma_0\), choose an \(\alpha\), and set an affinity/conductance resembling
\[
  c_{ij}^{\mathrm{P10}}
  =
  k_0\!\left(\frac{d_{ij}^2}{\sigma_0^2\widehat R_i\widehat R_j}\right)
  \frac{1}{\widehat R_i^\alpha\widehat R_j^\alpha},
\]
with the caveat that P10's theorem is written for samples on a smooth
manifold with Euclidean ambient distances. Translating to a graph or
Riemannian-complex setting is a derived design step, not a theorem already in
the paper.

The most conservative default for a smoothing comparator is \(\alpha=1\),
because it targets \(\Delta_p\) and \(E_p\), matching a density-weighted
smoothness interpretation. If the target is closer to geometry with uniform
volume, \(\alpha=1-d/2\) or the mixed \(\widehat\rho,\widehat p\)
normalization in the paper's Equation (21) becomes relevant, but then
intrinsic dimension or density estimation choices must be surfaced. For
\code{gflow}, \(\alpha\) should be an explicit parameter rather than hidden in
the conductance formula.

The strongest transfer point is the low-density variance comparison. If the
current overlap-density smoother and a P10-style adaptive kernel smoother
behave differently near sparse regions, P10 predicts why: the adaptive kernel
is designed to increase local averaging scale where sample density is low.
That is a reason to include the comparator. It is not evidence that the
overlap-density smoother is wrong, nor evidence that inverse-length
conductance and self-tuned kernel conductance are equivalent. They encode
different notions of local geometry.

\section{What To Reuse}

The review should reuse four elements from P10. First, use the normalized kNN
radius story to explain adaptive local scale: \(\widehat R_i\) estimates a
multiple of \(p(x_i)^{-1/d}\), so sparse regions get larger bandwidths.
Second, use \(\alpha\) to keep operator targets explicit. Third, use
Dirichlet-form convergence as the closest theoretical support for graph
smoothing penalties. Fourth, use the fixed-versus-self-tuned variance
comparison to motivate adaptive-kernel comparators under nonuniform sampling.

The review should avoid three overextensions. Do not cite P10 as a proof of
the current SIMODS overlap-density conductance. Do not describe the kNN radius
as a kNN graph-support theorem. Do not claim spectral convergence from the
MNIST or \(S^1\) figures. P10 is strongest as a mathematically careful bridge
from self-tuned spectral-clustering practice to continuum graph-Laplacian
limits.

\section*{Provenance}

Primary source: the local canonical P10 PDF. Archival scaffolding: the P10
Markdown memo and audit. Figure screenshots: internal-review captures listed
in the shared figure manifest, with included paths of the form
\code{../../paper\_figures/...}. This review follows the shared LaTeX style
file and keeps the current \code{fit.rdgraph.regression()} overlap-density
smoother separate from the planned P10 adaptive local-scale comparator.

\bibliographystyle{plainnat}
\bibliography{../../references}

\end{document}
